% Estilo compartido. Compilar el volumen principal con XeLaTeX o Tectonic.
\documentclass[12pt,letterpaper]{article}
\usepackage[spanish,es-noshorthands]{babel}
\usepackage[left=23mm,right=23mm,top=23mm,bottom=22mm,headheight=15pt]{geometry}
\usepackage{graphicx,xcolor,fancyhdr}
\usepackage[unicode,pdflang={es-MX},bookmarksopen,colorlinks=true,linkcolor=azul]{hyperref}
% La familia Latin Modern de 10 puntos también cubre los títulos grandes.
\DeclareFontShape{TU}{lmr}{m}{n}{<->\UnicodeFontFile{lmroman10-regular}{\UnicodeFontTeXLigatures}}{}
\DeclareFontShape{TU}{lmr}{b}{n}{<->\UnicodeFontFile{lmroman10-bold}{\UnicodeFontTeXLigatures}}{}
\DeclareFontShape{TU}{lmr}{bx}{n}{<->\UnicodeFontFile{lmroman10-bold}{\UnicodeFontTeXLigatures}}{}
\DeclareFontFamily{TU}{lmtt}{}
\DeclareFontShape{TU}{lmtt}{m}{n}{<->\UnicodeFontFile{lmmono10-regular}{}}{}
\definecolor{azul}{HTML}{075479}
\definecolor{tinta}{HTML}{172F3D}
\definecolor{muted}{HTML}{455E6A}
\definecolor{niebla}{HTML}{EDF4F6}
\renewcommand{\headrulewidth}{0pt}
\setlength{\parindent}{0pt}
\setlength{\parskip}{4pt}
\setlength{\emergencystretch}{1.3em}
\linespread{1.08}
\pagestyle{fancy}\fancyhf{}
\fancyhead[L]{\small\color{muted}ESCOM \enspace / \enspace Trabajo Terminal 2026-B155}
\fancyhead[R]{\small\color{muted}\volumen}
\fancyfoot[L]{\small\hyperlink{indice}{Volver al índice}}
\fancyfoot[R]{\small\color{muted}\thepage}
\newcommand{\boton}[1]{\textbf{#1}}
\newcommand{\tema}[1]{\par\addvspace{9pt}{\large\bfseries\color{azul}#1\par}\nobreak}
\newcommand{\pagina}[3]{\clearpage\phantomsection\addcontentsline{toc}{section}{#2}%
{\small\bfseries\color{azul}\MakeUppercase{#1}\par}\vspace{2pt}%
{\raggedright\fontsize{23}{26}\selectfont\bfseries\color{tinta}#2\par}\nobreak\vspace{2pt}%
{\large\color{muted}#3\par}\vspace{4pt}}
\newenvironment{pasos}{\begin{enumerate}\setlength{\itemsep}{3pt}\setlength{\topsep}{1pt}\setlength{\parskip}{0pt}}{\end{enumerate}}
\newcommand{\nota}[2]{\par\addvspace{7pt}\noindent\begingroup\setlength{\fboxsep}{8pt}%
\colorbox{niebla}{\parbox{\dimexpr\linewidth-2\fboxsep\relax}{\small\color{tinta}\textbf{#1}\enspace #2}}\endgroup\par}
\newcommand{\captura}[4][1]{\par\addvspace{6pt}\begin{center}\setlength{\fboxsep}{0pt}%
\fcolorbox{niebla}{white}{\includegraphics[width=\dimexpr#1\linewidth-2\fboxrule\relax,height=#3,keepaspectratio]{figuras/#2.png}}%
\par\vspace{2pt}{\small\color{muted}#4\par}\end{center}}
\newcommand{\portada}[2]{%
\hypersetup{pdftitle={Manual de usuario. #1},pdfauthor={Díaz Ramírez Daniel y Téllez De La Cruz Esaul},pdfsubject={Aplicación web para la gestión de actividades de la ESCOM. TT 2026-B155}}
\color{tinta}\thispagestyle{empty}
\begin{minipage}[c]{.18\linewidth}\includegraphics[width=.83\linewidth]{figuras/escudo-escom.png}\end{minipage}\hfill
\begin{minipage}[c]{.79\linewidth}{\large\bfseries Instituto Politécnico Nacional\par}Escuela Superior de Cómputo\par Trabajo Terminal 2026-B155\end{minipage}
\par\vspace{14mm}{\color{azul}\bfseries MANUAL DE USUARIO\par}\vspace{1mm}
{\fontsize{32}{35}\selectfont\bfseries #1\par}\vspace{2mm}
{\large Aplicación web para la gestión de actividades de la ESCOM\par}\vspace{4mm}
#2\par\vspace{4mm}
\hypertarget{indice}{}\tableofcontents
\vfill
{\small\color{muted}\textbf{Autores}\quad Díaz Ramírez Daniel y Téllez De La Cruz Esaul\par
\textbf{Directora}\quad Dra. Mónica Rivera De La Rosa\par
Edición de noviembre de 2026 \enspace · \enspace Ciudad de México\par}}
\AtBeginDocument{\renewcommand{\contentsname}{Índice de tareas}}
